% ============================================================
\chapter{Apprentissage Machine avec la TDA}
\label{ch:ml_tda}
% ============================================================

Les diagrammes de persistance capturent la ``forme'' des donn\'ees, mais les algorithmes d'apprentissage machine --- SVM, for\^ets al\'eatoires, r\'eseaux de neurones --- attendent des vecteurs de dimension fixe. Comment combler ce foss\'e~? C'est le probl\`eme de la \emph{vectorisation}, l'un des d\'efis centraux de la TDA appliqu\'ee. Peter Bubenik (2015) propose les \emph{paysages de persistance}, des fonctions dans un espace de Banach~; Henry Adams et ses collaborateurs introduisent les \emph{images de persistance}~; Mathieu Carri\`ere et al.\ d\'eveloppent les noyaux de persistance. Plus r\'ecemment, les couches de persistance diff\'erentiables permettent d'int\'egrer la TDA directement dans les r\'eseaux de neurones profonds, ouvrant la voie \`a un apprentissage topologique de bout en bout.

\begin{intuition}
Les diagrammes de persistance vivent dans un espace non vectoriel~:
on ne peut pas directement les additionner ni les multiplier par un
scalaire. Pour les utiliser dans un pipeline d'apprentissage machine,
il faut les \textbf{vectoriser}\index{vectorisation}, c'est-\`a-dire
les transformer en repr\'esentations vivant dans un espace de Hilbert
ou un espace de features de dimension finie.
\end{intuition}

% ============================================================
\section{Paysages de persistance}
\label{sec:landscapes}
% ============================================================

Rappelons le paysage de persistance introduit au
Chapitre~\ref{ch:diagrammes}~: pour un diagramme
$D = \{(b_i, d_i)\}$, le $k$-\`eme paysage $\lambda_k(t)$ est la
$k$-\`eme plus grande valeur des fonctions tentes $\Lambda_i(t)$.
Les paysages vivent dans un espace de Banach, ce qui permet de calculer
moyennes, variances et de r\'ealiser des tests statistiques.
\index{paysage de persistance}

\begin{minted}{python}
from gudhi.representations import Landscape
import numpy as np

# Diagramme de persistance
dgm = np.array([[0.0, 1.5], [0.5, 2.0], [1.0, 3.0]])

# Calcul des 3 premiers paysages sur une grille
landscape = Landscape(num_landscapes=3, resolution=100)
L = landscape.fit_transform([dgm])
print(f"Shape: {L.shape}")  # (1, 300)
\end{minted}

% ============================================================
\section{Images de persistance}
\label{sec:images}
% ============================================================

\begin{definition}[Image de persistance --- Adams et al.\ 2017]
\index{image de persistance}
Soit $D$ un diagramme de persistance. L'\textbf{image de persistance}
est construite en trois \'etapes~:
\begin{enumerate}
  \item \textbf{Changement de coordonn\'ees}~: on transforme chaque
        $(b_i, d_i)$ en $(b_i, d_i - b_i)$ (naissance, persistance).
  \item \textbf{Pond\'eration}~: on applique une fonction de poids
        $w : \R^2 \to \R_+$ (typiquement lin\'eaire en la persistance)
        pour att\'enuer les points proches de la diagonale.
  \item \textbf{Lissage gaussien}~: on place une gaussienne
        $\mathcal{N}((b_i, p_i), \sigma^2 I)$ sur chaque point et
        on discr\'etise sur une grille $N \times N$.
\end{enumerate}
L'image r\'esultante est~:
\[
  \rho(x, y) = \sum_{i=1}^n w(b_i, p_i) \cdot
  \frac{1}{2\pi\sigma^2}
  \exp\!\left(-\frac{(x - b_i)^2 + (y - p_i)^2}{2\sigma^2}\right).
\]
\end{definition}

\begin{theoreme}[Stabilit\'e des images de persistance]
Si $w$ est lipschitzienne et born\'ee, l'application
$D \mapsto \rho_D$ est lipschitzienne pour la distance de Wasserstein
$W_1$~:
\[
  \norm{\rho_{D_1} - \rho_{D_2}}_\infty \leq C \cdot W_1(D_1, D_2).
\]
\end{theoreme}

% ============================================================
\section{Silhouettes de persistance}
\label{sec:silhouettes}
% ============================================================

Rappelons la silhouette de persistance d'ordre $q$ introduite au
Chapitre~\ref{ch:diagrammes}~: c'est la moyenne pond\'er\'ee (par
$(d_i - b_i)^q$) des fonctions tentes. Pour $q$ grand, elle donne plus
de poids aux structures persistantes. La silhouette est un \'el\'ement
de $L^p(\R)$, donc directement utilisable comme feature dans un
pipeline d'apprentissage.
\index{silhouette de persistance}

% ============================================================
\section{M\'ethodes \`a noyaux}
\label{sec:noyaux}
% ============================================================

\begin{definition}[Noyau de persistance]
\index{noyau de persistance}
Un \textbf{noyau de persistance} est un noyau d\'efini positif
$K : \mathcal{D} \times \mathcal{D} \to \R$ sur l'espace des
diagrammes de persistance.
\end{definition}

\begin{definition}[Noyau \`a \'echelle de persistance --- Reininghaus et al.\ 2015]
\[
  k_\sigma(D_1, D_2) = \frac{1}{8\pi\sigma}
  \sum_{p \in D_1} \sum_{q \in D_2}
  \left(e^{-\frac{\norm{p - q}^2}{8\sigma}}
  - e^{-\frac{\norm{p - \bar{q}}^2}{8\sigma}}\right)
\]
o\`u $\bar{q} = (d_q, b_q)$ est le point $q$ r\'efl\'echi par rapport
\`a la diagonale.
\end{definition}

\begin{proposition}[D\'efini positivit\'e]
Le noyau $k_\sigma$ est d\'efini positif et stable~:
$\abs{k_\sigma(D_1, D_1') - k_\sigma(D_2, D_2')} \leq
C_\sigma (W_1(D_1, D_2) + W_1(D_1', D_2'))$.
\end{proposition}

\begin{definition}[Noyau de Fisher --- Le, Yamada 2018]
\index{noyau de Fisher}
Le \textbf{noyau de Fisher} mod\'elise chaque diagramme comme une
distribution et utilise la distance de Fisher entre distributions~:
\[
  K_{\mathrm{Fisher}}(D_1, D_2) = \exp\!\left(
  -\frac{1}{2t} d_{\mathrm{Fisher}}^2(D_1, D_2)\right).
\]
\end{definition}

% ============================================================
\section{Int\'egration dans un pipeline ML}
\label{sec:pipeline}
% ============================================================

\begin{exemple}[Pipeline complet de classification]
\begin{enumerate}
  \item \textbf{Donn\'ees}~: nuages de points ou s\'eries temporelles.
  \item \textbf{Filtration}~: Rips, Alpha ou sublevel set.
  \item \textbf{Persistance}~: calcul des diagrammes via \texttt{gudhi}
        ou \texttt{ripser}.
  \item \textbf{Vectorisation}~: images de persistance, paysages ou noyaux.
  \item \textbf{Mod\`ele}~: SVM, Random Forest ou r\'egression logistique.
\end{enumerate}
\end{exemple}

\begin{minted}{python}
from sklearn.pipeline import Pipeline
from sklearn.svm import SVC
from gudhi.representations import (
    PersistenceImage, Landscape, DiagramSelector
)

# Pipeline scikit-learn avec features topologiques
pipe = Pipeline([
    ("selector", DiagramSelector(use=True)),
    ("persistence_image", PersistenceImage(
        bandwidth=0.1, resolution=[20, 20]
    )),
    ("classifier", SVC(kernel="rbf", C=10.0)),
])

# pipe.fit(X_train_diagrams, y_train)
# y_pred = pipe.predict(X_test_diagrams)
\end{minted}

\begin{attention}[title=\textbf{Choix de la vectorisation}]
\begin{itemize}
  \item Les \textbf{paysages} pr\'eservent la structure de Banach et
        permettent des tests statistiques.
  \item Les \textbf{images} sont plus adapt\'ees aux CNN car elles
        produisent des repr\'esentations 2D.
  \item Les \textbf{noyaux} sont pr\'ef\'erables avec des SVM en
        petite dimension.
\end{itemize}
\end{attention}

% ============================================================
\section{Scikit-TDA et \'ecosyst\`eme logiciel}
\label{sec:scikit_tda}
% ============================================================

\begin{remarque}[\'Ecosyst\`eme Python pour la TDA]
\index{scikit-TDA}
\begin{itemize}
  \item \textbf{GUDHI}~: biblioth\`eque C++/Python compl\`ete
        (filtrations, persistance, repr\'esentations).
  \item \textbf{Ripser}~: calcul rapide de la persistance de Rips.
  \item \textbf{scikit-tda}~: m\'eta-paquet int\'egrant Ripser,
        Persim (distances), KeplerMapper.
  \item \textbf{giotto-tda}~: pipeline TDA compatible scikit-learn
        avec acc\'el\'eration C++.
\end{itemize}
\end{remarque}

% ============================================================
\section{Exercices}
% ============================================================

\begin{exercice}
Calculer les deux premiers paysages de persistance du diagramme
$D = \{(0, 3), (1, 4), (2, 5)\}$ et les tracer.
\end{exercice}

\begin{exercice}
Montrer que la silhouette de persistance d'ordre $q = 1$ est la
moyenne pond\'er\'ee (par la persistance) des fonctions tentes.
\end{exercice}

\begin{exercice}[Impl\'ementation]
Construire un pipeline complet qui classifie des nuages de points
\'echantillonn\'es sur le cercle vs.\ la figure en huit, en utilisant
des images de persistance et un SVM.
\end{exercice}

\begin{exercice}
Montrer que le noyau de Reininghaus est invariant par r\'eflexion
diagonale~: $k_\sigma(D, D') = k_\sigma(D', D)$.
\end{exercice}

\begin{exercice}[Comparaison exp\'erimentale]
Comparer les performances de classification (accuracy, temps de
calcul) entre paysages, images et noyaux de persistance sur un
jeu de donn\'ees de votre choix.
\end{exercice}

\begin{exercice}[Stabilit\'e des paysages de persistance]
Soient $D_1$ et $D_2$ deux diagrammes de persistance et
$\lambda_k^{(1)}$, $\lambda_k^{(2)}$ leurs $k$-\`emes paysages respectifs.
\begin{enumerate}
  \item Montrer que pour chaque point $(b, d)$ du diagramme, la fonction
        tente $\Lambda(t) = \max(0, \min(t - b, d - t))$ est $1$-lipschitzienne.
  \item En utilisant le fait que l'op\'eration $\mathrm{kmax}$ est
        $1$-lipschitzienne en norme $\ell^\infty$ sur les vecteurs tri\'es,
        montrer que~:
        \[
          \norm{\lambda_k^{(1)} - \lambda_k^{(2)}}_{L^\infty}
          \leq W_\infty(D_1, D_2).
        \]
  \item G\'en\'eraliser au cas $L^p$~: montrer que
        $\norm{\lambda_k^{(1)} - \lambda_k^{(2)}}_{L^p}
        \leq C_p \cdot W_p(D_1, D_2)^{1/p}$
        et discuter la constante $C_p$.
\end{enumerate}
\end{exercice}

\begin{exercice}[Pipeline SVM \`a noyau avec features de persistance]
On souhaite classifier des nuages de points en trois classes~:
cercle, tore (section 2D) et sph\`ere (\'echantillonn\'ee en 3D).
\begin{enumerate}
  \item Pour chaque nuage, calculer le diagramme de persistance en
        homologie de degr\'es 0 et 1 via un complexe de Vietoris--Rips.
  \item Construire la matrice de Gram du noyau de Reininghaus~:
        $G_{ij} = k_\sigma(D_i, D_j)$ pour un param\`etre $\sigma$ choisi
        par validation crois\'ee.
  \item Entra\^iner un SVM \`a noyau pr\'ecalcul\'e (\texttt{kernel='precomputed'})
        et rapporter l'accuracy en validation crois\'ee 5-folds.
  \item Comparer avec un pipeline utilisant des images de persistance
        concat\'en\'ees (degr\'es 0 et 1) et un SVM RBF. Discuter les
        avantages et inconv\'enients de chaque approche.
\end{enumerate}
\end{exercice}
